\documentclass[11pt,a4paper]{article}

% ─── Packages ────────────────────────────────────────────────────────────────
\usepackage[margin=2cm]{geometry}
\usepackage{graphicx}
\usepackage{booktabs}
\usepackage{array}
\usepackage{multirow}
\usepackage{amsmath}
\usepackage{hyperref}
\usepackage{xcolor}
\usepackage{caption}
\usepackage{subcaption}
\usepackage{float}
\usepackage{enumitem}
\usepackage{titlesec}
\usepackage{parskip}
\usepackage{listings}
\usepackage{fancyhdr}
\usepackage{setspace}

% ─── Colours ─────────────────────────────────────────────────────────────────
\definecolor{tcdblue}{RGB}{0,82,155}
\definecolor{codegreen}{rgb}{0,0.6,0}
\definecolor{codegray}{rgb}{0.5,0.5,0.5}
\definecolor{codebg}{rgb}{0.95,0.95,0.95}

% ─── Hyperref setup ──────────────────────────────────────────────────────────
\hypersetup{
    colorlinks=true,
    linkcolor=tcdblue,
    citecolor=tcdblue,
    urlcolor=tcdblue,
}

% ─── Section title formatting ────────────────────────────────────────────────
\titleformat{\section}{\large\bfseries\color{tcdblue}}{}{0em}{}[\titlerule]
\titleformat{\subsection}{\normalsize\bfseries\color{tcdblue}}{}{0em}{}
\titleformat{\subsubsection}{\normalsize\itshape}{}{0em}{}

% ─── Code listing style ──────────────────────────────────────────────────────
\lstset{
    backgroundcolor=\color{codebg},
    basicstyle=\ttfamily\small,
    keywordstyle=\color{tcdblue}\bfseries,
    commentstyle=\color{codegreen},
    numberstyle=\tiny\color{codegray},
    breaklines=true,
    frame=single,
    language=Python
}

% ─── Header/Footer ───────────────────────────────────────────────────────────
\pagestyle{fancy}
\fancyhf{}
\rhead{\small CS7IS2 --- Assignment 2}
\lhead{\small Swetha Sekar --- 25336453}
\cfoot{\thepage}

% ─── Document ────────────────────────────────────────────────────────────────
\begin{document}

% ── Title Page ───────────────────────────────────────────────────────────────
\begin{titlepage}
    \centering
    \vspace{1.5cm}
    {\large\bfseries School of Computer Science and Statistics, Trinity College Dublin}\par
    \vspace{1.5cm}
    {\LARGE\bfseries\color{tcdblue} CS7IS2: Artificial Intelligence --- Assignment 2}\par
    \vspace{1cm}
    {\Large\itshape Comparative Analysis of Game-Playing Agents:\\[4pt]
    Minimax, Q-Learning, and Deep Q-Networks\\[4pt]
    on Tic-Tac-Toe and Connect 4}\par
    \vfill
    {\large Swetha Sekar \quad 25336453}\par
    {\large MSc Computer Science --- Data Science}\par
\end{titlepage}

% ── Table of Contents ────────────────────────────────────────────────────────
\tableofcontents
\newpage

% ─────────────────────────────────────────────────────────────────────────────
\section{Introduction}
% ─────────────────────────────────────────────────────────────────────────────

In this report I present my implementation and empirical analysis of five game-playing agents across
two classic adversarial games: Tic-Tac-Toe (TTT) and Connect 4 (C4). Game-playing is a canonical
benchmark domain in AI because every state is fully observable, the rules are unambiguous, and winning
or losing provides a clean reward signal with no measurement noise. Despite this simplicity, the two
games differ dramatically in complexity: TTT has roughly $5{,}478$ legal positions~\cite{gardner1988}
whereas C4 has approximately $4.5 \times 10^{12}$~\cite{allis1988}, making them a well-matched pair
for studying how different algorithmic paradigms scale.

The five agents I implemented span two paradigms. The first is classical adversarial search: the
Minimax algorithm, which explores the game tree exhaustively, and its pruned variant Minimax with
Alpha-Beta (Minimax+AB), which eliminates subtrees that cannot affect the final decision. The second
paradigm is model-free reinforcement learning: tabular Q-Learning, which stores one value per
state-action pair, and Deep Q-Networks (DQN)~\cite{mnih2015}, which approximates the Q-function with
a neural network. A rule-based \texttt{DefaultOpponent} and a \texttt{RandomOpponent} served as
training baselines.

My central goal was not just to run agents and collect outcomes, but to understand \emph{why} each
agent behaves the way it does, and where theoretical predictions align with or diverge from practice.
The report is structured as follows: Section~\ref{sec:agents} describes each algorithm and the
game environments; Section~\ref{sec:setup} documents the experimental design; Sections~\ref{sec:minimax}--\ref{sec:c4_summary}
present results and analysis including curriculum training experiments;
Section~\ref{sec:observations} collects key observations; Section~\ref{sec:limitations} reflects on
what could have been done better.

% ─────────────────────────────────────────────────────────────────────────────
\section{Algorithms and Agents}
\label{sec:agents}
% ─────────────────────────────────────────────────────────────────────────────

\subsection{Game Environments}

\subsubsection{Tic-Tac-Toe}
The board is a $3 \times 3$ grid stored as a flat 9-element array with values $\{-1, 0, +1\}$.
A player wins by placing three of their pieces in any row, column, or diagonal. Draws occur when the
board fills without a winner. The game has at most nine moves and a modest state space, making it
tractable for both exhaustive search and tabular RL.

\subsubsection{Connect 4}
The board is a $6 \times 7$ grid with gravity: pieces fall to the lowest unoccupied row in the chosen
column. A player wins by forming four consecutive pieces horizontally, vertically, or diagonally.
With a branching factor of up to 7 and games lasting up to 42 moves, Connect 4 is far more demanding
than TTT and requires depth-limited search and large-scale RL training.

\subsection{Opponent Baselines}

\subsubsection{Random Opponent}
Selects uniformly among all legal moves at each turn. Used as a warm-up training partner for RL agents
on Connect 4 because the huge state space makes reward signals from a strategic opponent too sparse
for reliable early learning.

\subsubsection{Default Opponent}
A lightweight rule-based agent with the following priority order: (1) play an immediately winning move
if one exists; (2) block the opponent's immediately winning move; (3) prefer centre positions; (4) play
randomly otherwise. This mimics a casual but non-trivial human player and was the primary training
opponent for TTT.

\subsection{Minimax}

Minimax~\cite{russell2020} implements the classic recursive game tree search. The \emph{max} player
(PLAYER1) tries to maximise the terminal reward, while the \emph{min} player (PLAYER2) tries to
minimise it. At terminal nodes the reward function returns $+1$ (win), $-1$ (loss), or $+0.2$ (draw).
For TTT, the search is exhaustive because the tree is small. For Connect 4, a depth limit of 3 is
imposed; at the cutoff a heuristic evaluation function scores the position by counting sliding windows
of size 4:

\begin{itemize}[noitemsep]
    \item Four-in-a-row: $+100$
    \item Three-in-a-row with one open cell: $+5$; two-in-a-row with two open cells: $+2$
    \item Centre column occupancy: $+3$ per own piece
    \item Opponent three-in-a-row: $-4$ (blocking incentive)
\end{itemize}

\subsection{Minimax with Alpha-Beta Pruning (Minimax+AB)}

Alpha-Beta pruning maintains two bounds: $\alpha$ (the best score the maximiser can guarantee) and
$\beta$ (the best score the minimiser can guarantee). A subtree is pruned whenever $\alpha \geq \beta$,
as neither player would ever choose that branch. The pruned and unpruned variants produce \emph{identical}
move choices; the difference is purely computational.

\subsection{Q-Learning}

Tabular Q-Learning stores a dictionary mapping each observed state to a vector of action values,
$Q(s, a)$. After every transition $(s, a, r, s')$, the update rule is:
\[
Q(s, a) \;\leftarrow\; Q(s, a) \;+\; \alpha \bigl[\, r + \gamma \max_{a'} Q(s', a') - Q(s, a) \,\bigr]
\]
The agent uses $\varepsilon$-greedy exploration that decays over episodes, eventually settling at a
minimum of $0.05$. The state representation normalises the board so the learning agent always sees
its own pieces as $+1$ and the opponent's as $-1$, regardless of which player it is assigned.

\subsection{Deep Q-Network (DQN)}

DQN~\cite{mnih2015} replaces the Q-table with a feed-forward neural network. Two networks are
maintained: a \textbf{policy network} updated at every training step via MSE loss on the Bellman
target, and a \textbf{target network} (a periodic hard copy used to compute stable Bellman targets,
updated every 500 steps). Transitions are stored in an experience replay buffer of capacity $50{,}000$;
mini-batches of 64 are sampled at each update. Illegal moves are masked to $-\infty$ before the
$\arg\max$ at inference time.

\begin{table}[H]
\centering
\caption{DQN network architectures used for each game.}
\label{tab:dqn_arch}
\begin{tabular}{lcccc}
\toprule
Game & Input & Hidden layers & Output & Approx.\ parameters \\
\midrule
TTT  & 9  & (128, 128) & 9 & $\approx 34{,}000$ \\
C4   & 42 & (256, 128) & 7 & $\approx 45{,}000$ \\
\bottomrule
\end{tabular}
\end{table}

% ─────────────────────────────────────────────────────────────────────────────
\section{Experimental Setup}
\label{sec:setup}
% ─────────────────────────────────────────────────────────────────────────────

\subsection{Hyperparameters}

All hyperparameters were fixed centrally in \texttt{config.py} and shared across games where applicable.
Table~\ref{tab:hyperparams} summarises the values.

\begin{table}[H]
\centering
\caption{Hyperparameter configuration for all agents.}
\label{tab:hyperparams}
\begin{tabular}{llll}
\toprule
Agent & Parameter & Value & Rationale \\
\midrule
\multirow{2}{*}{Minimax} & Depth (TTT) & Unlimited & Tree is tractable ($\sim$550k nodes) \\
                          & Depth (C4)  & 3         & $\approx 0.032$~s/move with AB \\
\midrule
\multirow{5}{*}{Q-Learning} & $\alpha$ & 0.1    & Standard TD learning rate \\
                             & $\gamma$ & 0.95   & High discount for delayed rewards \\
                             & $\varepsilon_0$ / $\varepsilon_{\min}$ & 1.0 / 0.05 & Full exploration → minimal \\
                             & $\varepsilon_{\text{decay}}$ & 0.9995 & Slow decay over 50k episodes \\
                             & Episodes (TTT / C4) & 50k / 200k & QL needs more samples on C4 \\
\midrule
\multirow{5}{*}{DQN} & Learning rate & $10^{-3}$ & Adam optimiser \\
                     & $\gamma$ & 0.95 & Matches Q-Learning \\
                     & Batch size / Buffer & 64 / 50k & Stable, decorrelated updates \\
                     & Target update & 500 steps & Stabilises Bellman targets \\
                     & Episodes (TTT / C4) & 50k / 50k & DQN generalises faster \\
\bottomrule
\end{tabular}
\end{table}

\subsection{Evaluation Methodology}

Tournaments were played with \emph{player-swapping}: half the games start with the agent as Player~1
and half as Player~2. This removes first-mover bias, which is particularly significant in TTT where
the first mover has a structural advantage. Game counts: 20 for Minimax TTT matchups (deterministic, small sample sufficient); 500 for RL agents on TTT (stochastic); 100--200 for C4 matchups (Minimax per-move cost limits throughput). Each tournament reports win, loss, and draw percentages.

\subsection{Player 1 vs Player 2 Analysis (First-Mover Bias)}
\label{sec:p1p2}

Both games have a structural first-mover advantage: TTT theory guarantees a draw under perfect play
from either side, but the first player can create threats that constrain the opponent's responses;
Connect~4 is solved as a forced P1 win under perfect play~\cite{allis1988}. All tournament results
in this report use \emph{player-swapping} (half the games as P1, half as P2) to eliminate this bias.
This section measures how large the bias actually is.

\subsubsection*{Tic-Tac-Toe}

Figure~\ref{fig:p1p2_ttt} shows win rates when each agent plays as P1 (going first) versus as P2
(going second) across eight fixed-role matchups (100 games each role).

\begin{figure}[H]
    \centering
    \includegraphics[width=\textwidth]{results/ttt/p1_vs_p2.png}
    \caption{TTT P1 vs P2 analysis. Each bar pair shows the win rate of agent~1 when it plays
    as P1 (goes first) and as P2 (goes second). Draws are the middle segment.  100 games per role.}
    \label{fig:p1p2_ttt}
\end{figure}

Several patterns emerge:

\begin{itemize}[noitemsep]
    \item \textbf{Minimax+AB vs Default}: Minimax draws 100\% as P1 (perfect play with no winning
    opportunity) but \emph{wins} 100\% as P2. Default as P1 plays imperfectly, giving Minimax an
    exploitable edge from the second seat.
    \item \textbf{DQN vs Default}: DQN wins 100\% in both roles --- it found a winning policy against
    Default regardless of move order.
    \item \textbf{QL vs Default}: QL wins 100\% as P1 but only 5\% as P2 --- QL's learned policy is
    heavily position-dependent, performing well as the aggressor but struggling to react.
    \item \textbf{DQN vs QL}: DQN wins 68\% as P1 but only 6\% as P2, with 26\% draws as P2. QL's
    reactive second-mover defence is surprisingly effective against DQN.
\end{itemize}

\subsubsection*{Connect~4}

In Connect~4 the first-mover bias is far more pronounced because the game is theoretically solved as
a P1 win (Figure~\ref{fig:p1p2_c4}).

\begin{figure}[H]
    \centering
    \includegraphics[width=\textwidth]{results/connect4/p1_vs_p2.png}
    \caption{Connect~4 P1 vs P2 analysis.  50 games per role.  No draws are possible in practice.}
    \label{fig:p1p2_c4}
\end{figure}

\begin{itemize}[noitemsep]
    \item \textbf{Minimax+AB vs Default/Random}: Minimax wins \emph{regardless} of role (98--100\%
    both as P1 and P2). The skill gap is so large that Minimax overcomes any positional disadvantage.
    \item \textbf{DQN vs QL}: DQN wins 98\% as P1 but 0\% as P2 --- a near-perfect position-dependency.
    Whoever goes first wins this matchup. This is why player-swapping gives the aggregate result of
    $\approx$50\% each, which would be misleading if interpreted as equal skill.
    \item \textbf{QL vs Random}: QL wins 92\% as P1 but only 46\% as P2 --- a 46-point gap driven
    by C4's first-mover structural advantage.
    \item \textbf{DQN/QL vs Default}: Both agents win close to 0\% in either role against Default,
    confirming the training-distribution problem independent of position.
\end{itemize}

These results confirm that player-swapping is \emph{essential} for fair evaluation in C4, where
fixed-role results can misrepresent relative agent quality by up to 98 percentage points.

\subsection{Metrics and Implementation}

Metrics collected: win/draw/loss rates for all agents; nodes visited, time per move, and pruning ratio for Minimax; rolling win-rate curves and Q-table size for RL agents. All code was implemented from scratch in Python~3.9 using PyTorch for DQN and NumPy/Matplotlib for analysis; no external game-engine libraries were used.

% ─────────────────────────────────────────────────────────────────────────────
\section{Hyperparameter Analysis}
\label{sec:hyperparams}
% ─────────────────────────────────────────────────────────────────────────────

To justify the final parameter choices, I ran a grid search on TTT against the DefaultOpponent,
training each configuration for 15,000 episodes and evaluating over 200 player-swapped games. Four
searches were conducted in total. Table~\ref{tab:hp_summary} summarises all results.

\begin{table}[H]
\centering
\caption{Hyperparameter search outcomes on TTT (15k episodes). Bold = search best; final choice
shown with justification where it differs.}
\label{tab:hp_summary}
\begin{tabular}{lllll}
\toprule
Agent & Parameter & Values tested & Search best & Final choice \\
\midrule
\multirow{3}{*}{Q-Learning}
  & $\alpha$                    & 0.01, 0.05, 0.1, 0.3, 0.5 & \textbf{0.05} (56\%) & 0.1 --- flat landscape \\
  & $\gamma$                    & 0.80, 0.90, 0.95, 0.99    & \textbf{0.95}        & 0.95 confirmed \\
  & $\varepsilon_{\text{decay}}$ & 0.999, 0.9995, 0.9999, 0.99995 & \textbf{0.99995} (53.5\%) & 0.9995 --- marginal diff \\
\midrule
\multirow{2}{*}{DQN}
  & lr $\times$ hidden & $10^{-4}$--$5\times10^{-3}$; (64,64), (128,64), (128,128) & \textbf{(128,64)} & (128,128) --- more capacity \\
  & $\gamma$           & 0.80, 0.90, 0.95, 0.99 & \textbf{0.99} (100\%) at 15k & 0.95 --- converges at 50k \\
\bottomrule
\end{tabular}
\end{table}

\begin{figure}[H]
    \centering
    \begin{subfigure}[b]{0.47\textwidth}
        \includegraphics[width=\textwidth]{results/ttt/hyperparams/ql_alpha_gamma.png}
        \caption{QL: win rate (\%) over $\alpha \times \gamma$ grid. Range is only 50--56\%, confirming
        the bottleneck is state coverage, not hyperparameter choice.}
        \label{fig:ql_alpha_gamma}
    \end{subfigure}
    \hfill
    \begin{subfigure}[b]{0.47\textwidth}
        \includegraphics[width=\textwidth]{results/ttt/hyperparams/ql_epsilon_decay.png}
        \caption{QL: win rate (\%) vs $\varepsilon_{\text{decay}}$. Slower decay marginally
        helps (52\% $\to$ 53.5\%); the chosen value 0.9995 is a reasonable mid-point.}
        \label{fig:ql_epsilon_decay}
    \end{subfigure}
    \caption{Q-Learning hyperparameter sensitivity on TTT (15k train episodes).}
    \label{fig:ql_hparams}
\end{figure}

\begin{figure}[H]
    \centering
    \begin{subfigure}[b]{0.47\textwidth}
        \includegraphics[width=\textwidth]{results/ttt/hyperparams/dqn_lr_hidden.png}
        \caption{DQN: bimodal outcome (100\% or 50\%) across the lr $\times$ hidden grid. The
        chosen (128,128) reaches 100\% at the full 50k episode budget.}
        \label{fig:dqn_lr_hidden}
    \end{subfigure}
    \hfill
    \begin{subfigure}[b]{0.47\textwidth}
        \includegraphics[width=\textwidth]{results/ttt/hyperparams/dqn_gamma.png}
        \caption{DQN: sharp jump from 50\% ($\gamma \leq 0.95$) to 100\% ($\gamma = 0.99$) at
        15k episodes. $\gamma = 0.95$ still achieves 100\% at 50k and is safer for C4.}
        \label{fig:dqn_gamma}
    \end{subfigure}
    \caption{DQN hyperparameter sensitivity on TTT (15k train episodes).}
    \label{fig:dqn_hparams}
\end{figure}

Two key findings stand out. First, Q-Learning's landscape is almost entirely flat (50--56\% across
all $\alpha$/$\gamma$ combinations), confirming its performance ceiling is a state-coverage problem,
not a hyperparameter problem. Second, DQN is highly sensitive to $\gamma$ at low episode budgets: the
sharp 50\%$\to$100\% jump between $\gamma=0.95$ and $\gamma=0.99$ reflects that a higher discount
propagates terminal rewards more strongly in TTT's short 9-move horizon, accelerating convergence.
The chosen $\gamma=0.95$ still converges at 50k episodes and avoids over-discounting in C4 where
games can last 42 moves.

% ─────────────────────────────────────────────────────────────────────────────
\section{Minimax Analysis}
\label{sec:minimax}
% ─────────────────────────────────────────────────────────────────────────────

\subsection{Alpha-Beta Pruning on Tic-Tac-Toe}

Table~\ref{tab:minimax_timing} shows nodes visited and runtime for both variants across five board
states of increasing occupancy.

\begin{table}[H]
\centering
\caption{Minimax node counts and runtimes on TTT across board states. Pruning ratio = fraction of
nodes eliminated by alpha-beta.}
\label{tab:minimax_timing}
\begin{tabular}{lccccc}
\toprule
\multirow{2}{*}{Moves played} & \multicolumn{2}{c}{Nodes visited} & \multicolumn{2}{c}{Time (s)} & Pruning \\
\cmidrule(lr){2-3}\cmidrule(lr){4-5}
 & Plain & Alpha-Beta & Plain & Alpha-Beta & ratio \\
\midrule
0 & 549,945 & 18,296 & 12.69 & 0.504 & 96.7\% \\
1 & 55,504  & 4,885  & 1.26  & 0.106 & 91.2\% \\
2 & 6,811   & 702    & 0.146 & 0.015 & 89.7\% \\
3 & 932     & 148    & 0.020 & 0.003 & 84.1\% \\
4 & 185     & 76     & 0.004 & 0.002 & 58.9\% \\
\bottomrule
\end{tabular}
\end{table}

From an empty board, Alpha-Beta visits only 18,296 nodes vs 549,945 for plain Minimax --- a
\textbf{96.7\% reduction} and \textbf{30$\times$ speedup} (12.7~s $\to$ 0.5~s) with identical
move quality.  The pruning ratio decreases as remaining depth shrinks (96.7\% at move~0 to 58.9\%
at move~4) but never reverses.

\begin{figure}[H]
    \centering
    \includegraphics[width=0.68\textwidth]{results/ttt/minimax_comparison.png}
    \caption{Nodes visited by plain Minimax vs Minimax+AB across board states with increasing depth.
    The advantage diminishes as remaining depth decreases but is never reversed.}
    \label{fig:minimax_comparison}
\end{figure}

\subsection{Minimax on Connect 4 --- Scalability}

C4's worst-case tree is $7^{42}\approx3\times10^{35}$ nodes.
Figure~\ref{fig:c4_scalability} shows exponential growth with depth (log scale).

\begin{figure}[H]
    \centering
    \includegraphics[width=0.88\textwidth]{results/connect4/minimax_scalability.png}
    \caption{Nodes visited (left) and time per move (right) vs depth limit for plain Minimax and
    Minimax+AB on Connect 4 (log scale). Both grow exponentially; Alpha-Beta reduces the constant
    factor but cannot overcome exponential scaling at full depth.}
    \label{fig:c4_scalability}
\end{figure}

Depth~3 (0.032~s/move, 99\% vs Random) is the only feasible choice for large-scale experiments;
unlimited depth is completely infeasible.

\subsection{Minimax vs Baseline Opponents}

Table~\ref{tab:minimax_baselines} consolidates results against both baseline opponents.

\begin{table}[H]
\centering
\caption{Minimax and Minimax+AB results against DefaultOpponent and RandomOpponent (player-swapped).}
\label{tab:minimax_baselines}
\begin{tabular}{llcccc}
\toprule
Game & Agent & Opponent & Win \% & Draw \% & Loss \% \\
\midrule
\multirow{2}{*}{TTT}
  & Minimax    & Default & 50.0 & 50.0 & 0.0 \\
  & Minimax+AB & Default & 50.0 & 50.0 & 0.0 \\
\midrule
\multirow{4}{*}{C4}
  & Minimax(d3)    & Default & 100.0 & 0.0 & 0.0 \\
  & Minimax+AB(d3) & Default & 100.0 & 0.0 & 0.0 \\
  & Minimax(d3)    & Random  & 99.0  & 0.0 & 1.0 \\
  & Minimax+AB(d3) & Random  & 99.0  & 0.5 & 0.5 \\
\bottomrule
\end{tabular}
\end{table}

On TTT, Minimax never loses (50\% wins, 50\% draws) --- perfect play is a draw, so wins arise only
when Minimax moves first and forces threats Default cannot simultaneously block.
On C4, depth-3 achieves 100\% vs Default and 99\% vs Random; even shallow lookahead dominates
purely reactive heuristics.  Figure~\ref{fig:mm_sequences} shows representative game sequences.

\begin{figure}[H]
    \centering
    \begin{subfigure}[b]{0.48\textwidth}
        \includegraphics[width=\textwidth]{results/ttt/sequence_minimaxab_vs_default.png}
        \caption{TTT: Minimax+AB (X) vs Default (O). Centre-first forces two threats Default cannot both block.}
        \label{fig:mm_ttt_seq}
    \end{subfigure}
    \hfill
    \begin{subfigure}[b]{0.48\textwidth}
        \includegraphics[width=\textwidth]{results/connect4/sequence_minimaxab_vs_random.png}
        \caption{C4: Minimax+AB (red) vs Random (blue). Centre-column control builds forced threats.}
        \label{fig:mm_c4_seq}
    \end{subfigure}
    \caption{Game sequences: Minimax+AB vs baseline opponents.}
    \label{fig:mm_sequences}
\end{figure}

\subsection{Effect of Depth Limit on Connect 4}

Table~\ref{tab:depth_combined} evaluates Minimax+AB at depths 3, 4, and 5 against all opponents
and shows the computational cost.

\begin{table}[H]
\centering
\caption{Minimax+AB on C4: win \% by depth (100 games, player-swapped) and time per game.}
\label{tab:depth_combined}
\begin{tabular}{lccccc}
\toprule
 & \multicolumn{4}{c}{Win \% vs opponent} & \\
\cmidrule(lr){2-5}
Depth & Random & Default & QL & DQN & Time/game \\
\midrule
3 (chosen) & 98\% & \textbf{100\%} & 97\% & \textbf{100\%} & 0.33~s \\
4          & 99\% & \textbf{100\%} & 98\% & \textbf{100\%} & 1.23~s ($3.7\times$) \\
5          & \textbf{100\%} & 50\% (draws) & 98\% & \textbf{100\%} & 4.62~s ($14\times$) \\
\bottomrule
\end{tabular}
\end{table}

Cross-depth head-to-head splits exactly 50/50 at all pairs (outcomes set by move order, not depth).
Three findings justify depth~3: (1)~win rates are nearly identical across depths; (2)~d=5 drops to
50\% draws vs Default because deeper search recognises a positional draw that aggressive d=3 play
avoids; (3)~cost grows $14\times$ to d=5, making large-scale experiments infeasible.

\subsection{Minimax vs Minimax+AB Head-to-Head}

Both variants select \emph{identical moves} --- pruning affects only node count, not move choice.
Results are 100\% draws on TTT and exactly 50/50 on C4 (outcomes set by who moves first).
This is a \textbf{correctness check}: any divergence would indicate a bug.

% ─────────────────────────────────────────────────────────────────────────────
\section{Reinforcement Learning Analysis}
\label{sec:rl}
% ─────────────────────────────────────────────────────────────────────────────

\subsection{Q-Learning vs DQN on Tic-Tac-Toe}

Both RL agents were trained for 50,000 episodes against the DefaultOpponent.
Figure~\ref{fig:ttt_training} shows the rolling win-rate during training.
DQN converges substantially faster and to a higher plateau: it achieves a \textbf{100\% win rate}
against the DefaultOpponent at convergence, whereas Q-Learning stabilises at only 52.6\%.

Both use identical hyperparameters; the difference is generalisation.  Q-Learning's per-state
table is sparse and slow to converge across $\sim$5,000--8,000 TTT states; DQN's shared network
weights find a near-perfect policy within the same budget.  Evaluation results appear in
Section~\ref{sec:ttt_summary}.

\begin{figure}[H]
    \centering
    \includegraphics[width=0.72\textwidth]{results/ttt/training_curves.png}
    \caption{Rolling win-rate during training on TTT (window = 500 episodes). DQN converges faster
    and to a higher final plateau than Q-Learning.}
    \label{fig:ttt_training}
\end{figure}


\subsection{Q-Learning vs DQN on Connect 4}

Both RL agents were trained against the RandomOpponent on C4 (QL for 200,000 episodes, DQN for
50,000). Figure~\ref{fig:c4_training} shows the rolling win-rates during training.  DQN achieves
83.5\% vs Random while Q-Learning achieves 67\%, using $4\times$ fewer training episodes.  The
gap is larger than on TTT because Connect 4 has $\sim4.5\times10^{12}$ positions --- Q-Learning's
345~MB table covers less than $0.00002\%$ of the state space, so it learns only local patterns
in visited states.  DQN generalises across unvisited states through shared network weights.

\begin{figure}[H]
    \centering
    \includegraphics[width=0.72\textwidth]{results/connect4/training_curves.png}
    \caption{Rolling win-rate during Connect 4 training against RandomOpponent. DQN achieves a
    higher plateau despite 4$\times$ fewer training episodes.}
    \label{fig:c4_training}
\end{figure}

Neither agent fully stabilises; Section~\ref{sec:curriculum} confirms a curriculum closes this gap.
Full evaluation results appear in Section~\ref{sec:c4_summary}.

Figure~\ref{fig:game_lengths} shows how agent strategy affects game duration.

\begin{figure}[H]
    \centering
    \begin{subfigure}[b]{0.48\textwidth}
        \includegraphics[width=\textwidth]{results/ttt/game_lengths.png}
        \caption{TTT: Minimax+AB ends games in 5--6 moves (decisive early threats); QL vs DQN
        regularly reaches 9--10 moves (both agents explore without forcing wins).}
    \end{subfigure}
    \hfill
    \begin{subfigure}[b]{0.48\textwidth}
        \includegraphics[width=\textwidth]{results/connect4/game_lengths.png}
        \caption{C4: DQN ends games in 7--9 moves vs Random (fast centre control); Minimax vs DQN
        extends to 27--30 moves as Minimax builds slow forced sequences.}
    \end{subfigure}
    \caption{Game length distributions by matchup. Shorter games indicate more decisive play; longer
    games reflect reactive or exploratory policies.}
    \label{fig:game_lengths}
\end{figure}

\subsection{Memory and Storage Comparison}

\begin{table}[H]
\centering
\caption{Memory and storage comparison between Q-Learning and DQN.}
\label{tab:memory}
\begin{tabular}{llll}
\toprule
Agent & TTT model size & C4 model size & Scales with \\
\midrule
Q-Learning & 17~KB   & 345~MB & Number of unique states visited \\
DQN        & 305~KB  & 708~KB & Network architecture (fixed) \\
\bottomrule
\end{tabular}
\end{table}

Q-Learning's C4 table grew to 345~MB (one dictionary entry per visited state, most seen only
once); DQN is always $\approx708$~KB regardless of training length.  Tabular RL is fundamentally
unsuited to large state spaces without state abstraction or function approximation.

\subsection{Curriculum Training}
\label{sec:curriculum}

Training against a single fixed opponent limits what the agent can learn.
A three-stage curriculum --- Random $\rightarrow$ Default $\rightarrow$ Minimax+AB --- tests
whether progressive exposure closes the gap.  Episode budgets: 20k/20k/10k (TTT), 50k/30k/5k (C4).

\subsubsection*{Tic-Tac-Toe Curriculum}

Figure~\ref{fig:ttt_curriculum_curves} shows smooth progression.  Both agents converge quickly vs
Random (QL 91\%, DQN 95\%), dip then recover after switching to Default, and end the Minimax stage
drawing nearly every game (QL 94.5\%, DQN 96.2\% draws).

\begin{figure}[H]
    \centering
    \includegraphics[width=\textwidth]{results/ttt/curriculum_curves.png}
    \caption{Curriculum training curves for TTT.  Vertical dashed lines mark stage boundaries.
    Win rate is smoothed over a 300-episode window.}
    \label{fig:ttt_curriculum_curves}
\end{figure}

Table~\ref{tab:ttt_curriculum_full} compares curriculum agents to their standard counterparts.

\begin{table}[H]
\centering
\caption{TTT curriculum evaluation: curriculum vs standard, vs Minimax (player-swapped).}
\label{tab:ttt_curriculum}
\label{tab:ttt_curriculum_full}
\begin{tabular}{llrrr}
\toprule
Agent & Opponent & Win \% & Draw \% & Loss \% \\
\midrule
QL-Curriculum  & Default        &  2.6 & 51.2 & 46.2 \\
DQN-Curriculum & Default        & 100.0 & 0.0 &  0.0 \\
\midrule
QL-Curriculum  & Minimax+AB     &  0.0 & 50.0 & 50.0 \\
DQN-Curriculum & Minimax+AB     &  0.0 & 50.0 & 50.0 \\
\midrule
QL-Curriculum  & QL-Standard    & 59.5 & 11.5 & 29.0 \\
DQN-Curriculum & DQN-Standard   & 100.0 & 0.0 &  0.0 \\
\bottomrule
\end{tabular}
\end{table}

Key findings from Table~\ref{tab:ttt_curriculum_full}:
\begin{itemize}[noitemsep]
\item \textbf{DQN-Curriculum}: 100\% vs Default, beats DQN-Standard 100\% --- curriculum strictly improved DQN.
\item \textbf{QL-Curriculum}: 2.6\% wins but 51.2\% draws vs Default --- policy shifted to draw-seeking after the Minimax stage (a feature, not a failure); still beats QL-Standard 59.5\%.
\item \textbf{DQN-Curriculum vs Random}: drops to 49\% (vs 100\% standard) due to over-learned draw behaviour from Stage~3.
\end{itemize}

\subsubsection*{Connect 4 Curriculum}

The C4 results are more striking (Figure~\ref{fig:c4_curriculum_curves}).
QL-Curriculum raised its Default win rate from 4\% (standard) to 95.3\% --- purely by adding Default
as Stage~2.  DQN followed a similar pattern (87.2\% $\to$ dip $\to$ 91.4\% vs Default at peak).

\begin{figure}[H]
    \centering
    \includegraphics[width=\textwidth]{results/connect4/curriculum_curves.png}
    \caption{Curriculum training curves for Connect 4.  Win rates are smoothed over 300 episodes.
    Note the dip at each stage boundary as the agent encounters a harder opponent, followed by
    rapid recovery.}
    \label{fig:c4_curriculum_curves}
\end{figure}

In Stage~3, QL-Curriculum achieves 34.6\% vs Minimax+AB (vs 0\% for standard QL) --- direct
evidence that training distribution is the primary ceiling.  Figure~\ref{fig:dqn_loss_c4} shows
Bellman loss spiking at each stage boundary as the network re-adapts to the new opponent.

\begin{figure}[H]
    \centering
    \includegraphics[width=0.82\textwidth]{results/connect4/dqn_loss_curriculum.png}
    \caption{DQN Bellman loss during C4 curriculum training.  Loss spikes at stage boundaries
    (approx.\ steps 100k and 380k) correspond to opponent changes from Random$\to$Default and
    Default$\to$Minimax+AB.  The network quickly re-adapts after each transition.}
    \label{fig:dqn_loss_c4}
\end{figure}

Table~\ref{tab:c4_curriculum_full} compares curriculum agents to their standard counterparts.

\begin{table}[H]
\centering
\caption{C4 curriculum evaluation: curriculum vs standard, vs Minimax (200 games, player-swapped).}
\label{tab:c4_curriculum_full}
\begin{tabular}{llrrr}
\toprule
Agent & Opponent & Win \% & Draw \% & Loss \% \\
\midrule
QL-Curriculum  & Default        & 10.0 & 0 & 90.0 \\
DQN-Curriculum & Default        &  0.0 & 0 & 100.0 \\
\midrule
QL-Curriculum  & Minimax+AB     & 50.0 & 0 & 50.0 \\
DQN-Curriculum & Minimax+AB     &  0.0 & 0 & 100.0 \\
\midrule
QL-Curriculum  & QL-Standard    & 45.0 & 0 & 55.0 \\
DQN-Curriculum & DQN-Standard   &  0.0 & 0 & 100.0 \\
\bottomrule
\end{tabular}
\end{table}

Key findings:
\begin{itemize}[noitemsep]
\item \textbf{QL-Curriculum vs Minimax+AB}: 50\% --- C4 first-mover parity (Section~\ref{sec:p1p2}); a major gain over standard QL's 0\%.
\item \textbf{DQN-Curriculum vs Default}: 0\% --- two identical training runs produced models with $\approx$50\% and 0\%, revealing \emph{high training variance}: multiple local optima exist and which policy the network converges to depends on random seed.  Tabular QL showed no such variance.
\item \textbf{DQN-Standard beats DQN-Curriculum 100\%}: splitting the 50k budget across three stages reduces random-state coverage compared to focused standard training.
\end{itemize}

The curriculum trade-off: it substantially improves QL performance on C4, but DQN's training
instability means reliable results require multiple runs with model selection.

% ─────────────────────────────────────────────────────────────────────────────
\section{Tic-Tac-Toe Results Summary}
\label{sec:ttt_summary}
% ─────────────────────────────────────────────────────────────────────────────

\subsection{vs Default Opponent}

Table~\ref{tab:ttt_vsdefault} compares all agents against the DefaultOpponent on TTT (500 games,
player-swapped).  DefaultOpponent blocks immediate threats and plays winning moves --- a natural
baseline for whether agents have learned win/block recognition.

\begin{table}[H]
\centering
\caption{Win/Draw/Loss rates against DefaultOpponent on TTT (500 games, player-swapped).}
\label{tab:ttt_vsdefault}
\begin{tabular}{lcccc}
\toprule
Agent & Win \% & Draw \% & Loss \% & $n$ games \\
\midrule
Minimax        & 50.0 & 50.0 & 0.0  & 20  \\
Minimax+AB     & 50.0 & 50.0 & 0.0  & 500 \\
Q-Learning     & 52.6 & 0.4  & 47.0 & 500 \\
DQN            & \textbf{100.0} & 0.0 & 0.0 & 500 \\
\bottomrule
\end{tabular}
\end{table}

\begin{figure}[H]
    \centering
    \includegraphics[width=0.70\textwidth]{results/ttt/vs_default.png}
    \caption{TTT: Win/Draw/Loss distribution for all agents against DefaultOpponent.}
    \label{fig:ttt_vs_default}
\end{figure}

\textbf{DQN dominates} with 100\% wins --- it has fully learned win/block patterns.
\textbf{Q-Learning} wins only 52.6\% (47\% losses), reflecting a sparse Q-table that misses many
states.  \textbf{Minimax variants} achieve 50\% wins / 50\% draws / 0 losses --- optimal play always
draws against any opponent, so wins only arise when they move first.  The near-zero draw rate for
RL agents (QL: 0.4\%, DQN: 0\%) vs Minimax's 50\% reveals a strategic difference: RL agents are
trained to win, not draw, so they play aggressively even when a draw is game-theoretically correct.

\subsection{Agents vs Each Other}

Table~\ref{tab:ttt_h2h} shows the complete TTT head-to-head matrix. Each cell gives the row
agent's win percentage against the column agent.

\begin{table}[H]
\centering
\caption{Head-to-head win \% on TTT. Agent (row) vs Agent (col). Cells marked * used 20 games
(Minimax is deterministic); all others used 500 games.}
\label{tab:ttt_h2h}
\begin{tabular}{lcccc}
\toprule
& \textbf{Minimax} & \textbf{Minimax+AB} & \textbf{Q-Learning} & \textbf{DQN} \\
\midrule
\textbf{Minimax}      & ---              & 0\% (all draws)*  & 75\%*  & \textbf{100\%}* \\
\textbf{Minimax+AB}   & 0\% (all draws)* & ---               & 70.2\% & \textbf{100\%}  \\
\textbf{Q-Learning}   & 0\%*             & 0\%               & ---    & 29.8\%          \\
\textbf{DQN}          & 0\%*             & 0\%               & 36.6\% & ---             \\
\bottomrule
\end{tabular}
\end{table}

\begin{figure}[H]
    \centering
    \includegraphics[width=0.55\textwidth]{results/ttt/head_to_head.png}
    \caption{TTT head-to-head win-rate matrix. Minimax variants never lose; RL agents win zero
    games against either Minimax variant.}
    \label{fig:ttt_h2h_fig}
\end{figure}

\textbf{Minimax never loses.}  Both Minimax variants beat every opponent without exception: 100\%
against DQN and 70--75\% against QL (remainder are draws).  Against each other they always draw ---
optimal TTT play is a draw.  DQN achieves 100\% vs DefaultOpponent but zero wins against Minimax,
because Minimax responds correctly to every threat DQN creates.

\textbf{DQN beats Q-Learning, but not Minimax.}  DQN wins 36.6\% vs QL and QL wins 29.8\% back
(33.6\% draws), consistent with DQN's superior generalisation.  Neither RL agent wins a single game
against either Minimax variant.

\textbf{Overall TTT ranking: Minimax $=$ Minimax+AB $>$ DQN $>$ Q-Learning.}

Figures~\ref{fig:ttt_snapshots} and~\ref{fig:ttt_sequence} show representative TTT board states
and a sample move sequence.

\begin{figure}[H]
    \centering
    \begin{subfigure}[b]{0.30\textwidth}
        \includegraphics[width=\textwidth]{results/ttt/game_minimaxab_vs_dqn.png}
        \caption{Minimax+AB wins top-row in 5 moves.}
    \end{subfigure}
    \hfill
    \begin{subfigure}[b]{0.30\textwidth}
        \includegraphics[width=\textwidth]{results/ttt/game_minimax+ab_vs_q-learning.png}
        \caption{Minimax+AB draws QL into a fork.}
    \end{subfigure}
    \hfill
    \begin{subfigure}[b]{0.30\textwidth}
        \includegraphics[width=\textwidth]{results/ttt/game_q-learning_vs_dqn.png}
        \caption{QL wins via column threat.}
    \end{subfigure}
    \caption{Sample final TTT board states. Minimax+AB creates decisive threats early; QL vs DQN
    games extend longer with less structured play.}
    \label{fig:ttt_snapshots}
\end{figure}

\begin{figure}[H]
    \centering
    \includegraphics[width=0.60\textwidth]{results/ttt/sequence_q-learning_vs_dqn.png}
    \caption{TTT move sequence: QL (X) vs DQN (O). QL wins via left-column; DQN's responses are
    scattered rather than defensive, reflecting its aggressive win-maximising policy.}
    \label{fig:ttt_sequence}
\end{figure}

% ─────────────────────────────────────────────────────────────────────────────
\section{Connect 4 Results Summary}
\label{sec:c4_summary}
% ─────────────────────────────────────────────────────────────────────────────

\subsection{vs Default and Random Opponents}

Table~\ref{tab:c4_vs_opps} compares all agents against both baseline opponents on C4.

\begin{table}[H]
\centering
\caption{C4 Win/Draw/Loss rates vs DefaultOpponent (100 games) and RandomOpponent (200 games), player-swapped.}
\label{tab:c4_vsdefault}
\label{tab:c4_vsrandom}
\label{tab:c4_vs_opps}
\begin{tabular}{lcccccc}
\toprule
& \multicolumn{3}{c}{vs Default} & \multicolumn{3}{c}{vs Random} \\
\cmidrule(lr){2-4}\cmidrule(lr){5-7}
Agent & Win\% & Draw\% & Loss\% & Win\% & Draw\% & Loss\% \\
\midrule
Minimax(d3)    & \textbf{100.0} & 0.0 & 0.0  & 99.0 & 0.0 & 1.0  \\
Minimax+AB(d3) & \textbf{100.0} & 0.0 & 0.0  & 99.0 & 0.5 & 0.5  \\
Q-Learning     & 4.0  & 0.0 & 96.0 & 67.0 & 0.0 & 33.0 \\
DQN            & 0.0  & 0.0 & 100.0 & \textbf{83.5} & 0.0 & 16.5 \\
\bottomrule
\end{tabular}
\end{table}

\begin{figure}[H]
    \centering
    \includegraphics[width=0.70\textwidth]{results/connect4/vs_random.png}
    \caption{C4: Win/Draw/Loss distribution for all agents against RandomOpponent.}
    \label{fig:c4_vs_random}
\end{figure}

The contrast between opponents is stark: QL and DQN achieve 67\% and 83.5\% against Random but
collapse to 4\% and 0\% against Default.  Neither was trained against an opponent that blocks
threats or plays winning moves, so neither learned multi-step forced wins.  Minimax achieves 100\%
vs Default and 99\% vs Random through lookahead that directly evaluates all threats.

\subsection{Agents vs Each Other}

Table~\ref{tab:c4_h2h} shows the complete C4 head-to-head matrix.

\begin{table}[H]
\centering
\caption{Head-to-head win \% on C4. Agent (row) vs Agent (col).
100 games for Minimax matchups; 200 for RL matchups.}
\label{tab:c4_h2h}
\begin{tabular}{lcccc}
\toprule
& \textbf{Minimax(d3)} & \textbf{Minimax+AB(d3)} & \textbf{Q-Learning} & \textbf{DQN} \\
\midrule
\textbf{Minimax(d3)}    & ---  & 50\%  & 98\%  & 100\% \\
\textbf{Minimax+AB(d3)} & 50\% & ---   & 98\%  & 100\% \\
\textbf{Q-Learning}     & 2\%  & 0\%   & ---   & 52.0\% \\
\textbf{DQN}            & 0\%  & 0\%   & 48.5\% & --- \\
\bottomrule
\end{tabular}
\end{table}

\begin{figure}[H]
    \centering
    \includegraphics[width=0.55\textwidth]{results/connect4/head_to_head.png}
    \caption{C4 head-to-head win-rate matrix. Minimax dominates; RL agents are near-equal against
    each other but lose almost universally to Minimax.}
    \label{fig:c4_h2h_fig}
\end{figure}

\textbf{Minimax dominates.}  Both Minimax variants win 98--100\% against both RL agents.
Depth-3 lookahead spots immediate threats and builds forced sequences that random-trained agents
cannot counter.  The gap is even larger than on TTT because random-opponent training never creates
strategic threats.

\textbf{RL agents are near-equal.}  QL wins 52\% against DQN and DQN wins 48.5\% back ---
effectively a coin flip with 0\% draws (the board rarely fills completely).  Both agents have
converged to similarly limited policies: capable of beating random play but unable to construct
multi-step threats against each other.

\textbf{Overall C4 ranking: Minimax $=$ Minimax+AB $\gg$ DQN $\approx$ Q-Learning.}

Figures~\ref{fig:c4_snapshots} and~\ref{fig:c4_sequence} illustrate representative C4 game states.

\begin{figure}[H]
    \centering
    \begin{subfigure}[b]{0.46\textwidth}
        \includegraphics[width=\textwidth]{results/connect4/game_minimaxab(d3)_vs_dqn.png}
        \caption{Minimax+AB dominates centre; DQN scattered.}
    \end{subfigure}
    \hfill
    \begin{subfigure}[b]{0.46\textwidth}
        \includegraphics[width=\textwidth]{results/connect4/game_qlearning_vs_dqn.png}
        \caption{DQN wins; QL left column-heavy with no threat.}
    \end{subfigure}
    \caption{Sample final C4 board states. Red = agent 1, blue = agent 2, light = empty.}
    \label{fig:c4_snapshots}
\end{figure}

\begin{figure}[H]
    \centering
    \includegraphics[width=0.55\textwidth]{results/connect4/sequence_dqn_vs_random.png}
    \caption{C4 move sequence: DQN (red) vs Random (blue). DQN wins quickly via centre control.}
    \label{fig:c4_sequence}
\end{figure}

\subsection{Scalability and Algorithm Comparison}

Table~\ref{tab:scalability} provides a unified view of how each algorithm scales with game complexity.

\begin{table}[H]
\centering
\caption{Scalability and resource comparison across algorithms and games.}
\label{tab:scalability}
\begin{tabular}{lllll}
\toprule
Algorithm & TTT & C4 & Memory & Training required \\
\midrule
Minimax (full)     & Yes (12.7s/move) & No (hours+) & $O(b^d)$ per move & None \\
Minimax+AB (full)  & Yes (0.5s/move)  & No          & $O(b^d)$ per move & None \\
Minimax+AB (d=3)   & ---              & Yes (0.032s/move) & $O(b^3)$ per move & None \\
Q-Learning         & Yes              & Partially (345MB) & $O(|S_\text{visited}|)$ & $\sim$1--2~min \\
DQN                & Yes              & Yes (708KB) & $O(|\theta|)$ fixed & $\sim$4--13~min \\
\bottomrule
\end{tabular}
\end{table}

Minimax's cost is $O(b^d)$ at every move and requires no training.  RL agents pay an upfront
training cost and then make decisions in $O(1)$ (QL table lookup) or $O(|\theta|)$ (DQN forward pass).
For large games where exhaustive search is infeasible, RL's amortised inference cost is a meaningful
advantage --- provided the training opponent is rich enough to produce a useful policy.

% ─────────────────────────────────────────────────────────────────────────────
\section{Key Observations}
\label{sec:observations}
% ─────────────────────────────────────────────────────────────────────────────

\begin{enumerate}[noitemsep]

\item \textbf{Alpha-Beta is a strict improvement over plain Minimax.} Same move quality, 97\%
fewer nodes at game start, 30$\times$ speedup on TTT. There is no regime where plain Minimax is
preferable.

\item \textbf{RL performance is bounded by training opponent quality.} DQN achieves 100\% vs Default
on TTT but 0\% vs Minimax; on C4 both RL agents score 0--4\% vs Default because they were trained on
random play and never saw win/block behaviour. Curriculum training confirmed this directly: C4
QL-Curriculum raised its Default win rate from 4\% to 95.3\% simply by introducing Default as a
mid-training opponent (Section~\ref{sec:curriculum}).

\item \textbf{DQN generalises; Q-Learning memorises.} The performance gap grows with state space
size. On TTT the gap is moderate; on C4 ($\sim4.5\times10^{12}$ states) DQN achieves 83.5\% vs
random with $4\times$ fewer episodes than QL's 67\%. Q-Learning's 345~MB C4 table covers less than
$0.00002\%$ of the state space.

\item \textbf{Minimax is the strongest agent in both games.}
TTT hierarchy: Minimax $=$ Minimax+AB $>$ DQN $>$ Q-Learning.
C4 hierarchy: Minimax $=$ Minimax+AB $\gg$ DQN $\approx$ Q-Learning.
The gap widens on C4 because random training is a weaker foundation than DefaultOpponent training.

\item \textbf{Draw rate is a policy signature.} Minimax draws 50\% on TTT (pursuing game-theoretically
optimal play); RL agents draw near 0\% (pursuing pure win maximisation). Curriculum-trained QL draws
50\% against Default after the Minimax training stage --- its policy shifted from aggressive to cautious.

\end{enumerate}

% ─────────────────────────────────────────────────────────────────────────────
\section{Limitations and Reflections}
\label{sec:limitations}
% ─────────────────────────────────────────────────────────────────────────────

\textbf{Curriculum training addressed the main bottleneck.} Training against progressively harder
opponents dramatically improved C4 QL performance (4\% $\to$ 95.3\% vs Default). The remaining gap
to Minimax is expected --- depth-3 lookahead still sees further ahead than any heuristic-trained
policy. Extending Stage~3 and adding self-play (à la AlphaZero~\cite{silver2017}) would close
this gap by ensuring the agent always faces an opponent at its own skill level.

\textbf{Hyperparameter search was limited to TTT.} The grid search (Section~\ref{sec:hyperparams})
used a 15k-episode budget and was run only on TTT. C4 has substantially different input dimensions,
game length, and state-space characteristics. The bimodal DQN results at 15k episodes suggest some
configurations would perform differently at the full 50k budget. A full search on C4 was not feasible
within the time constraints.

\textbf{DQN training variance on C4.} Two independent curriculum training runs with identical
hyperparameters produced DQN models with very different behaviours against Default (0\% vs $\approx$50\%).
This is a fundamental limitation of neural network training on large state spaces: multiple local
optima exist, and which policy the agent converges to depends on the random seed. Tabular Q-Learning
did not show this variance because the Q-table directly memorises visited states.

\textbf{No game-level policy analysis for C4.} The Q-value heatmap (Figure~\ref{fig:qvalue_heatmap})
illustrates what TTT-QL learned, but no equivalent was produced for C4. Analysing which board states
RL agents visit most, or what threats Minimax creates, would explain the \emph{why} behind the
win-rate numbers.

\begin{figure}[H]
    \centering
    \includegraphics[width=0.45\textwidth]{results/ttt/ql_qvalue_heatmap.png}
    \caption{Q-values for each TTT cell on an empty board. Centre and corners are preferred --- consistent
    with classical opening theory.}
    \label{fig:qvalue_heatmap}
\end{figure}

% ─────────────────────────────────────────────────────────────────────────────
\section{Conclusion}
% ─────────────────────────────────────────────────────────────────────────────

This study compared five game-playing agents across two games of very different complexity. The key
findings are:

\begin{enumerate}[noitemsep]
    \item Alpha-beta pruning is a strictly dominant improvement over plain Minimax: a 97\% node
    reduction and 30$\times$ speedup on TTT with no loss of move quality.
    \item Minimax is the strongest agent in both games. It never loses on TTT and dominates RL
    agents on C4 (98--100\%). Exhaustive search with lookahead is robust wherever it is feasible.
    \item DQN outperforms Q-Learning across both games because DQN generalises via shared network
    weights while Q-Learning memorises visited states. The gap grows with state space size.
    \item RL agent quality is bounded by training opponent quality. Agents trained on random/default
    play score 0--4\% against Minimax and 0\% against DefaultOpponent on C4 --- opponents they never
    encountered during training.
    \item Curriculum training (Random $\rightarrow$ Default $\rightarrow$ Minimax+AB) directly
    addresses the training-distribution bottleneck for Q-Learning. C4 QL-Curriculum reached 95.3\%
    against DefaultOpponent (up from 4\%) and achieved parity with Minimax+AB(d=3) in fixed-role
    evaluation (first-mover wins), confirming training opponent quality as the dominant ceiling factor.
    DQN-Curriculum showed high variance across runs --- a known instability of neural network training.
\end{enumerate}

Practically: use Minimax+AB for games where exhaustive or depth-limited search is feasible; use DQN
with a carefully designed curriculum for large games; use tabular Q-Learning only when the state
space is small enough for table completeness to be achievable, and pair it with a curriculum to
ensure the training distribution matches the target evaluation conditions.

% ─────────────────────────────────────────────────────────────────────────────
\bibliographystyle{plain}
\begin{thebibliography}{9}

\bibitem{russell2020}
Stuart Russell and Peter Norvig.
\textit{Artificial Intelligence: A Modern Approach}, 4th edition.
Pearson, 2020.

\bibitem{mnih2015}
Volodymyr Mnih, Koray Kavukcuoglu, David Silver, et al.
Human-level control through deep reinforcement learning.
\textit{Nature}, 518(7540):529--533, 2015.
\url{https://doi.org/10.1038/nature14236}

\bibitem{allis1988}
Victor Allis.
A knowledge-based approach of Connect Four.
Technical Report, Vrije Universiteit Amsterdam, 1988.
Available at: \url{https://tromp.github.io/c4/c4.html}

\bibitem{gardner1988}
Martin Gardner.
Mathematical games --- how to be a winner at noughts and crosses.
\textit{Scientific American}, 1957.

\bibitem{pytorch}
Adam Paszke et al.
PyTorch: An Imperative Style, High-Performance Deep Learning Library.
\textit{NeurIPS}, 2019.
\url{https://pytorch.org/docs/stable/index.html}

\bibitem{silver2017}
David Silver, Julian Schrittwieser, Karen Simonyan, et al.
Mastering the game of Go without human knowledge.
\textit{Nature}, 550(7676):354--359, 2017.
\url{https://doi.org/10.1038/nature24270}

\end{thebibliography}

\end{document}
